% ============================================================
%  VS Code Light+ 多语言代码高亮
%  支持：C / C++ / Python / Java / Verilog / VHDL / ASM
% ============================================================

% --- 颜色定义（模拟 VS Code Light+ 主题）---
\definecolor{vscBg}{RGB}{255,255,255}
\definecolor{vscComment}{RGB}{0,128,0}
\definecolor{vscKeyword}{RGB}{0,0,255}
\definecolor{vscString}{RGB}{163,21,21}
\definecolor{vscNumber}{RGB}{9,134,136}
\definecolor{vscPreproc}{RGB}{0,0,255}
\definecolor{vscType}{RGB}{38,127,153}
\definecolor{vscFunc}{RGB}{121,94,38}
\definecolor{vscLineNum}{RGB}{127,127,127}
\definecolor{vscBorder}{RGB}{210,210,210}
\definecolor{vscLeftBar}{RGB}{0,122,204}     % 左侧竖线色

% ============================================================
%  基础样式（带细线框）
% ============================================================
\lstdefinestyle{vscode-base}{
    backgroundcolor=\color{vscBg},
    basicstyle=\footnotesize\ttfamily,
    columns=flexible,
    keepspaces=true,
    extendedchars=false,
    breaklines=true,
    breakatwhitespace=false,
    tabsize=4,
    showtabs=false,
    showspaces=false,
    commentstyle=\color{vscComment}\itshape,
    keywordstyle=\color{vscKeyword}\bfseries,
    stringstyle=\color{vscString},
    numberstyle=\tiny\color{vscLineNum},
    emphstyle=\color{vscType},
    numbers=left,
    numbersep=8pt,
    frame=single,
    rulecolor=\color{vscBorder},
    frameround=tttt,
    xleftmargin=1.5em,
    framexleftmargin=1.5em,
    aboveskip=1.0em,
    belowskip=1.0em,
    captionpos=b,
    abovecaptionskip=0.5em,
    belowcaptionskip=0.8em,
    % 行号背景
    numberblanklines=true,
    escapeinside={@*}{*@},                   % 允许在代码块内插入 LaTeX：@*\LaTeX{}*@
}

% ============================================================
%  简洁样式（仅左侧竖线，无框 — 适合短代码）
% ============================================================
\lstdefinestyle{vscode-simple}{
    style=vscode-base,
    frame=leftline,
    framesep=8pt,
    framerule=3pt,
    rulecolor=\color{vscLeftBar},
    xleftmargin=0pt,
    framexleftmargin=0pt,
    numbers=none,
}

% ============================================================
%  C / Embedded C（含嵌入式关键字）
% ============================================================
\lstdefinelanguage{EmbeddedC}[]{C}{
    morekeywords={uint8_t,uint16_t,uint32_t,uint64_t,
        int8_t,int16_t,int32_t,int64_t,
        size_t,ssize_t,bool,true,false,NULL,
        volatile,inline,register,restrict,_Bool,
        Xil_In8,Xil_Out8,Xil_In16,Xil_Out16,
        Xil_In32,Xil_Out32,
        XPAR_AXI_GPIO_0_BASEADDR,XPAR_AXI_GPIO_1_BASEADDR,
        XPAR_AXI_TIMER_0_BASEADDR},
    sensitive=true,
}

\lstdefinestyle{vscode-c}{
    style=vscode-base,
    language=EmbeddedC,
    morecomment=[l]{//},
    morecomment=[s]{/*}{*/},
    emph={int,void,char,short,long,float,double,
        unsigned,signed,struct,union,enum,
        typedef,const,static,extern,return,
        if,else,for,while,do,switch,case,break,continue,
        goto,sizeof},
}

\lstdefinestyle{vscode-c-pp}{
    style=vscode-c,
    directivestyle=\color{vscPreproc},
}

% ============================================================
%  C++
% ============================================================
\lstdefinestyle{vscode-cpp}{
    style=vscode-base,
    language=C++,
    morekeywords={nullptr,constexpr,noexcept,override,final,
        typename,template,namespace,using,class,public,private,protected,
        virtual,explicit,static_cast,dynamic_cast,unique_ptr,shared_ptr},
}

% ============================================================
%  Python
% ============================================================
\lstdefinestyle{vscode-python}{
    style=vscode-base,
    language=Python,
    emph={self,cls,True,False,None},
    morekeywords={as,async,await,nonlocal,match,case,with,lambda,yield,raise,except,finally},
}

% ============================================================
%  Java
% ============================================================
\lstdefinestyle{vscode-java}{
    style=vscode-base,
    language=Java,
    morekeywords={extends,implements,interface,abstract,
        synchronized,transient,volatile,instanceof,
        package,import,super},
}

% ============================================================
%  Verilog / SystemVerilog（FPGA 开发必备）
% ============================================================
\lstdefinelanguage{Verilog}{
    morekeywords={
        module,endmodule,
        input,output,inout,
        wire,reg,tri,wand,wor,
        supply0,supply1,
        assign,always,initial,
        posedge,negedge,or,
        begin,end,fork,join,
        if,else,case,endcase,casex,casez,default,
        for,while,repeat,forever,
        parameter,localparam,defparam,
        integer,real,time,realtime,
        signed,unsigned,
        function,endfunction,
        task,endtask,
        generate,endgenerate,genvar,
        specify,endspecify,
        `define,`ifdef,`else,`endif,`include,`timescale,
    },
    morecomment=[l]{//},
    morecomment=[s]{/*}{*/},
    morestring=[b]",
    sensitive=true,
}

\lstdefinestyle{vscode-verilog}{
    style=vscode-base,
    language=Verilog,
    emph={input,output,inout,wire,reg,parameter,localparam,integer},
}

% SystemVerilog 扩展
\lstdefinelanguage{SystemVerilog}[Verilog]{Verilog}{
    morekeywords={
        logic,bit,byte,shortint,int,longint,
        always_comb,always_ff,always_latch,
        enum,struct,union,typedef,
        interface,endinterface,
        modport,clocking,
        package,endpackage,
        rand,randc,constraint,
        assert,assume,cover,
        new,this,super,
        virtual,extends,implements,
    },
}

\lstdefinestyle{vscode-systemverilog}{
    style=vscode-base,
    language=SystemVerilog,
}

% ============================================================
%  VHDL（FPGA 开发备选）
% ============================================================
\lstdefinelanguage{VHDL}{
    morekeywords={
        entity,end,architecture,of,is,begin,
        port,in,out,inout,buffer,linkage,
        signal,variable,constant,type,subtype,
        process,component,generic,map,
        if,then,elsif,else,end,case,when,
        for,loop,while,next,exit,
        wait,on,until,report,severity,
        assert,after,others,
        std_logic,std_logic_vector,
        bit,bit_vector,boolean,integer,real,
        unsigned,signed,std_ulogic,
        rising_edge,falling_edge,
        library,use,all,package,body,
        function,procedure,return,
        downto,to,range,
        and,or,nand,nor,xor,xnor,not,
        IEEE,STD_LOGIC_1164,NUMERIC_STD,
    },
    morecomment=[l]{--},
    morestring=[b]",
    sensitive=false,
}

\lstdefinestyle{vscode-vhdl}{
    style=vscode-base,
    language=VHDL,
    emph={entity,architecture,process,component,port,generic,map,
        signal,variable,constant,std_logic,std_logic_vector,
        IEEE,STD_LOGIC_1164,NUMERIC_STD},
}

% ============================================================
%  x86 / ARM 汇编
% ============================================================
\lstdefinelanguage{x86asm}{
    morekeywords={
        mov,push,pop,add,sub,mul,div,inc,dec,
        and,or,xor,not,shl,shr,rol,ror,
        cmp,test,jmp,je,jne,jz,jnz,jg,jl,jge,jle,
        call,ret,int,iret,nop,hlt,
        db,dw,dd,dq,equ,org,section,global,extern,
        eax,ebx,ecx,edx,esi,edi,ebp,esp,
        ax,bx,cx,dx,si,di,bp,sp,
        al,ah,bl,bh,cl,ch,dl,dh,
        rax,rbx,rcx,rdx,rsi,rdi,rbp,rsp,r8,r9,r10,r11,r12,r13,r14,r15,
    },
    morecomment=[l]{;},
    morestring=[b]",
    sensitive=false,
}

\lstdefinestyle{vscode-asm}{
    style=vscode-base,
    language=x86asm,
    emph={eax,ebx,ecx,edx,esi,edi,ebp,esp,eip,rax,rbx,rcx,rdx},
}

% ============================================================
%  默认样式
% ============================================================
\lstset{style=vscode-base}

% ============================================================
%  行内代码
% ============================================================
\lstdefinestyle{vscode-inline}{
    basicstyle=\small\ttfamily,
    commentstyle=\color{vscComment},
    keywordstyle=\color{vscKeyword},
    stringstyle=\color{vscString},
    breaklines=false,
    showstringspaces=false,
    backgroundcolor=\color{vscBg},
}

% ============================================================
%  引用外部源文件
% ============================================================
\newcommand{\includecode}[2][]{%
    \lstinputlisting[style=vscode-base,#1]{#2}%
}

% 按语言语法糖
\newcommand{\includeC}[2][]{%
    \lstinputlisting[style=vscode-c-pp,#1]{#2}%
}
\newcommand{\includePython}[2][]{%
    \lstinputlisting[style=vscode-python,#1]{#2}%
}
\newcommand{\includeVerilog}[2][]{%
    \lstinputlisting[style=vscode-verilog,#1]{#2}%
}
\newcommand{\includeVHDL}[2][]{%
    \lstinputlisting[style=vscode-vhdl,#1]{#2}%
}
\newcommand{\includeASM}[2][]{%
    \lstinputlisting[style=vscode-asm,#1]{#2}%
}
